\section{Version control}
\label{sec:version_control}
Let one recall the Modularity (M) and the Development (D) requirements. In essence, the architecture has to be modular in the sense that it has to enable the development of software modules that can be easily maintained as standalone units, but at the same time it should provide a straightforward integration process when a larger-scale project, like a prototype, requires to set up many of these modules at the same time. Furthermore, when working in any of these two situations, the architecture should provide easy means of maintaining the single units composing a larger project and of tracking the changes made to them while working on the project itself. Moreover, it should be possible to issue updates to these modules and propagate them to other projects that integrate the same modules, in case they do not require a specific version of them. The technology that addresses these requirements is called \emph{version control}.

Version control \cite{zolkifli_version_2018} is an essential element of contemporary software design and development, serving as a systematic mechanism for managing changes to source code, documentation, and other artifacts associated with a project over time. Version control involves the practice of tracking and managing modifications to software artifacts, thereby allowing developers to maintain a comprehensive history of changes, collaborate efficiently, and ensure consistent software quality. Version control systems (VCSs) enable concurrent development by multiple contributors without conflicts, thus mitigating issues such as overwritten code and conflicting updates. Furthermore, VCS facilitate the exploration of new features or bug fixes through branching, which isolates development efforts into independent streams, and merging, which integrates these branches into a cohesive codebase. This adaptability is crucial for modern software teams, who have to concurrently manage multiple development lines, such as feature development, bug resolution, and experimental changes.

The criticality of version control in software development is profound. It provides detailed tracking of every individual change made by team members, which facilitates an understanding of software evolution and captures the rationale behind each modification. This meticulous tracking becomes invaluable during debugging, as developers can analyze the precise changes that contributed to a malfunction and revert to a previous stable version if necessary. In large-scale software projects, where multiple developers are working simultaneously on a variety of features, a clear history of changes helps ensure accountability and transparency, both of which are vital for quality assurance and effective collaboration. The ability to revert changes or establish a timeline of modifications ensures that teams can efficiently correct mistakes without sacrificing valuable work or disrupting ongoing development processes. Ultimately, version control transcends the role of a mere tool for tracking changes—it is a foundational practice that underlies efficient, collaborative, and high-quality software development.

For these reasons, it is clear that the adoption of a version control system as a basic element of the architecture provides the means to satisfy the (M) and (D) requirements: as for the Modularity (M) requirement, the version control system enables to maintain the software ranging from modules to larger projects as standalone repositories under version control, eventually linked among each other. Furthermore, thanks to the branching and merging capabilities of the version control system, the development process is streamlined, and the integration of the modules into a larger project is facilitated. The version control system also provides the means to track the changes made to the software, to issue updates to the modules, and to propagate them to other projects that integrate them, thus satisfying the Development (D) requirement. The inner workings of these processes are described in the following section, which introduces the VCS of choice for the proposed architecture.

\subsection{Git}
\label{sec:git}
Git \cite{torvalds_git_nodate} is one of the most influential version control systems in contemporary software development \cite{loeliger_version_2012,bird_promises_2009}, recognized for its robustness, efficiency, and distributed architecture. Its inception dates back to 2005, when Linus Torvalds, the creator of the Linux operating system \cite{torvalds_linux_nodate}, sought an effective and flexible version control system for managing the Linux kernel's development. At that time, the Linux kernel project was experiencing significant growth, with thousands of developers from around the globe contributing code. Initially, the project relied on a proprietary version control system called BitKeeper, which, although effective, became unsuitable for the open-source nature of the kernel due to licensing conflicts. This prompted Torvalds to develop Git, a system specifically tailored to meet the unique requirements of the kernel project, particularly distributed collaboration, speed, and support for non-linear workflows. There is no official explanation of its name, but some well-known theories suggest that it may be a play on the word ``get'' or some underground slang term in an unidentified language meaning ``simple'' or ``stupid'', also since Torvalds himself refers to it as \emph{the stupid content tracker} \cite{torvalds_git_nodate}.
\begin{figure}[t!]
  \centering
  \includesvg[width=.35\textwidth]{images/ch3/git_logo.svg}
  \caption{Git logo \cite{torvalds_git_nodate}.}
  \label{fig:gitlogo}
\end{figure}

Linus Torvalds designed Git with several essential requirements in mind, derived from the practical challenges faced during the Linux kernel's rapid and distributed development. First, it needed to accommodate a distributed development model, whereby each developer maintains a complete copy of the entire codebase, including its full history. This decentralized approach provided resilience, ensuring that developers were not dependent on a central repository and could work independently without risking the loss of progress. The distributed nature also meant that individual contributors could experiment freely, with the assurance that any inadvertent changes do not compromise the project's integrity. The requirement for speed was equally paramount, as the system needed to handle the scale of contributions typical of the Linux project efficiently and support quick retrieval of project states. Furthermore, Torvalds emphasized the importance of strong support for non-linear development workflows, a feature critical for managing the extensive branching, merging, and contributions that characterize an open-source project of such magnitude and diversity.

The distributed architecture of Git is one of its defining strengths, distinguishing it from centralized version control systems \cite{rodriguez-bustos_how_2012}. In a centralized model, all developers rely on a single repository, which serves as the central point of authority for the project. In contrast, Git's distributed model allows each contributor to maintain a complete local clone of the repository, which includes the full history of all changes. This design not only provides a safety net, as the loss of a single repository does not impact the project's integrity, but also enhances productivity by allowing developers to work offline and synchronize changes when network access becomes available. In open-source projects like the Linux kernel, where developers are distributed across different geographies and time zones, the distributed model provides autonomy, enabling contributors to make progress independently while maintaining synchronization with the broader team. The distributed architecture also facilitates peer-to-peer collaboration, allowing developers to exchange changes directly without relying on a central authority, thereby promoting a more resilient and scalable development workflow.

Git's architecture and features make it a uniquely powerful and flexible version control system, which has revolutionized how teams approach software development. One significant advantage of Git is its use of snapshots rather than deltas. Unlike earlier version control systems, which typically stored differences between successive versions (deltas), Git creates a snapshot of the complete project state at each commit. This design choice ensures that the project's state can be easily and efficiently retrieved at any point in history, and it enhances data integrity through the use of cryptographic hashing. Each snapshot is uniquely identified using a hash function, which not only guarantees the integrity of the entire repository but also makes the history robust against corruption. This snapshot model forms the basis for the entire Git system, making retrieval, comparison, and merging of different states exceptionally efficient.

Git's branching and merging capabilities are also notable strengths, allowing developers to create branches with minimal overhead and merge them seamlessly when ready. Branches in Git are lightweight, which encourages experimentation without destabilizing the main codebase. This capability is particularly valuable for projects that require multiple, concurrent lines of development, such as new feature development, bug resolution, or exploratory work. The ability to create branches easily means that developers can work on different features or tasks simultaneously without fear of conflicting changes, which is vital in large-scale collaborative environments. Git's merging tools are designed to handle complex merging scenarios effectively, providing features such as conflict markers and automatic merging, which significantly reduce the complexity and risks associated with combining different branches of development.

Another advanced feature that Git offers, particularly beneficial in large and modular projects, is the concept of submodules. Submodules allow a Git repository to include and track the content of other repositories, effectively enabling complex projects to integrate external dependencies while maintaining version control over them. Submodules are indispensable in managing dependencies, particularly when different components of a larger system are developed independently. By incorporating submodules, developers can ensure that specific versions of external libraries or components are consistently tracked, which guarantees compatibility and stability across different versions of the main project. This is especially useful when projects rely on third-party libraries that need to remain stable across multiple updates, preventing unforeseen compatibility issues that may arise from updates that are not synchronized.

Submodules also play a key role in enabling modular development practices. Large projects often comprise multiple components that are logically separate but need to be version-controlled together. Submodules provide a means to maintain these components as distinct units, each with its own repository, while ensuring that they can be brought together in a cohesive manner within the parent project. This modularity is invaluable when different teams are responsible for different components, as each team can work independently on their respective modules, while the main project maintains coherence and consistency. Submodules thus foster a more manageable and scalable development process, particularly for projects that involve multiple dependencies with distinct development cycles. By leveraging submodules, teams can maintain a balance between modular autonomy and integrated consistency, ensuring that all parts of the system work harmoniously together while evolving independently.

The use of submodules extends beyond dependency management to foster more effective collaboration across different teams. When multiple teams work on separate but interrelated components, submodules provide an effective mechanism for versioning each component while maintaining an overarching structure. Each submodule is linked to a particular commit in its respective repository, which allows the main project to keep track of stable versions of dependencies. If updates are needed, the submodule can be updated to point to a new commit, ensuring that any changes are deliberate and controlled. This controlled mechanism ensures that all updates are tested before being integrated, which prevents disruptions to the main project's stability. Furthermore, submodules make it easier to reuse code across different projects, as they provide a clear and organized way to incorporate existing repositories without duplicating code, thereby fostering a more efficient and reusable development ecosystem.

A similar and complementary features to submodules is that of subtrees. Subtrees and submodules are both mechanisms for managing dependencies in a Git repository, but they offer distinct workflows and advantages. Subtrees allow developers to embed the content of an external repository into a subdirectory of the main repository. Unlike submodules, subtrees do not require additional configuration or extra repositories to be manually initialized when cloning the main repository. This makes subtrees more straightforward for new collaborators, as the embedded content becomes a natural part of the repository's history and structure, without any extra setup steps.

On the other hand, submodules are more suited for maintaining clear separations between distinct modules or components that need to evolve independently while remaining linked to a specific version in the main project. Submodules provide a finer level of control and facilitate updates across multiple repositories, which is particularly useful in larger projects where different teams are responsible for maintaining different components. Thus, while subtrees offer simplicity and convenience by integrating dependencies directly into the repository, submodules provide flexibility and modularity, enabling a cleaner separation of concerns.

These foundational features have led to Git becoming the de facto standard for version control in the software industry, trusted for both small-scale projects and large, highly distributed software systems. Git's flexibility, efficiency, and robust support for collaborative development have established it as an indispensable tool for modern software engineering, facilitating convoluted workflows and enabling distributed teams to innovate with confidence. Its influence extends beyond individual projects, shaping the entire software development landscape by enabling more effective collaboration, version management, and deployment practices. As software systems grow increasingly complex and teams become more distributed, the role of Git in managing the intricacies of software evolution remains pivotal, making it an essential element of any serious software development endeavor, as the one that this work proposes. To this end, the design of the architecture starts from the organization of every work it supports as a Git repository, as illustrated in the following chapters. To fully satisfy the requirements (M) and (D), the architecture also includes support for both submodules and subtrees, as they are both useful in different contexts and for different purposes. While submodules are suited to allow for the combined integration and development of single software units into larger projects, they require a more complex setup and management, which is not always necessary. Subtrees, on the other hand, are more straightforward to use and are more suited in scenarios in which it is not expected that active development on the embedded content is carried out, but rather that it is used as a dependency and eventually only updated to newer versions. The architecture is designed to support both mechanisms, allowing for the best of both worlds to be exploited.
